\chapter{AI discussion}

The usage of the AI in the framework of this report focuses on three pillars.

\begin{enumerate}
    \item \textbf{Intuition behind the mathematical developments : }The major usage of the AI was centered around an assistant role for decrypting unknown mathematical simplification or decomposition.  The mathematical fondations regarding the channel physics or information theory are well known for the electronics engineers.  Nevertheless, some details like the Fano's inequality or the Bessel function of first order was just slightly brushed in our course without a deep interrogation about the meaning.  The AI in this case could give us detailed undertanding and lead us to the intuition hidden behind the mathematics.
    \item \textbf{Help about the report construction : }The second idea was to use the AI as support tool for the direction to give to this report.  More specifically, it was to confirm the result or analysis-oriented report with a focus on the intuition behind the mathematical expressions instead of just summarize the research paper.
    \item \textbf{Support for example construction and validation of results : } Finally, AI was used as a tool to construct simplified numerical examples providing support in the discussion of the scenarii results or for the critical analysis. For instance, the example about the key generation time was supported by the artificial intelligence to confirm the number of OFDM symbols contained in a 5GNR communication packet and double-checked in the litterature afterwards.
\end{enumerate}



